\begin{table}[H]
\centering
\caption{Effect of the Deutschlandticket on Regional Labor Markets}
\label{tab:main}
\begin{threeparttable}
\begin{adjustbox}{max width=\textwidth}
\begin{tabular}{lccccc}
\toprule
 & (1) & (2) & (3) & (4) & (5) \\
 & Unemp. & Unemp. & Unemp. & Empl. & Unemp. \\
\midrule
Subsidy $\times$ Post & $0.0099$ & $0.0987$ & $0.1813$ & $0.1597$ & $0.0987$ \\
 & (0.0125) & (0.1254) & (0.2899) & (0.1267) & (0.0950) \\
 \\
\midrule
Treatment unit & EUR & EUR 10 & High/Low & EUR 10 & EUR 10 \\
Dep. variable & Unemp. & Unemp. & Unemp. & Empl. & Unemp. \\
Region FE & Yes & Yes & Yes & Yes & Yes \\
Year FE & Yes & Yes & Yes & Yes & Yes \\
Clustering & NUTS1 & NUTS1 & NUTS1 & NUTS1 & NUTS2 \\
Observations & 553 & 553 & 553 & 570 & 553 \\
Regions & 38 & 38 & 38 & 38 & 38 \\
\bottomrule
\end{tabular}
\end{adjustbox}
\begin{tablenotes}[flushleft]
\small
\item \textit{Notes:} Standard errors clustered at the NUTS1 level (16 states) in parentheses, except column~(5) which clusters at NUTS2 (38 regions). * $p<0.10$, ** $p<0.05$, *** $p<0.01$. The dependent variable is the annual regional unemployment rate (\%, columns 1--3, 5) or employment rate (\%, column 4). ``Subsidy'' is $\max(0, \text{pre-reform pass price} - 49)$. Column~(1) uses the raw EUR subsidy; (2) normalizes per EUR~10; (3) uses a binary indicator for above-median subsidy (EUR~23.5); (4) uses the employment rate as outcome; (5) uses NUTS2-level clustering. Data: Eurostat lfst\_r\_lfu3rt and lfst\_r\_lfe2emprt, 2010--2024.
\end{tablenotes}
\end{threeparttable}
\end{table}
